\quellseite{212}
\part[Part V. Helpful Hints]{Part V. Helpful Hints on the Five Parts of the Catechism}

\quellseite{214}
\chapter{Chapter XX. The Ten Commandments}

\randnotiz{Definition.} 1. Give a clear definition of God’s Law, e.\,g., it is the expression of God’s will.

God has one will. The Ten Commandments show ten different ways of offending against the one will. Question on each point.

\randnotiz{Historic sketch.} 2. Give brief historic sketch of God’s law-giving, before Sinai, at Sinai, after the Ten Commandments. References through Old Testament. New Testament. Prepare carefully. Refer to your Old and New Testament Theology, e.\,g., show law before fall, after fall, to Noah, the patriarchs, etc. Then the giving of the Ten Commandments to Moses, and by him to the people, etc. Question on each point.

\randnotiz{Kinds of law.} 3. Define clearly and distinguish clearly between (a) natural law, i. e., the laws of nature and the witness of conscience; (b) civil law; (c) ceremonial law; (d) moral. Illustrate and question on each point.

\randnotiz{Three uses of moral law.} 4. Define clearly, so that children can understand and distinguish clearly between the three uses of the moral law, viz., the political, the pedagogical and the didactic. Illustrate and question on each point.
\quellseite{215}
Show how the three uses are going to apply in the study of the Commandments by the class. Make clear the distinction between a legalistic spirit of obedience and an evangelical spirit, viz., that the former tries to obey because he must, the latter because he wants to, and that this latter spirit of obedience is the one that the baptized children of God should pray for and strive after continually; show how one might yield a seeming outward obedience, in act and in observing forms of worship, and fail to have the inner spirit of obedience. Illustrate and question.

\randnotiz{Law for man’s good.} 5. Show that God gives each Commandment for man’s good in this life and in the next; that He introduces the Ten Commandments by reminding man how kind He has been and is, and that thus He desires to impress His children that this law is for their good in body and in soul, in this world and in the next. Illustrate from each Commandment. Exemplify and question on each point.

\randnotiz{Relation of two tables.} 6. Show the relation of the two tables to each other, and how that no one can love his neighbor as he should unless he loves his God supremely. Distinguish between a mere humanitarian and a Christian love, between the narrow and selfish love taught and shown in the secret lodge and the broad and unselfish love of the Christian. Illustrate. Good Samaritan. Other examples. Question on each point.

\quellseite{216}
\randnotiz{Relation to each other.} 7. Show relation of the Commandments to each other, that all express the one will of God, e.\,g., in the first, God wants an undivided heart, satisfied with and loyal to the one God. In the second He wants a loyal and consecrated tongue. In the third He wants the whole man, body and soul, to have rest and refreshment, and to render the joyful homage and worship of heart and tongue.

In the second table God begins with those who stand next to Him, our parents and superiors, and demands the proper respect and obedience for our own well-being, for the comfort of the superiors, and for the glory of His name. In the fifth God would conserve the life and welfare of the body; in the sixth, the home and family; in the seventh, property; in the eighth, reputation and character, and in the ninth and tenth the purity of the inner springs of the soul.

After thus explaining the meaning and bearing of the moral law as a whole, you are ready to proceed to the separate Commandments.

\section{On the First Commandment}

1. Do not be afraid of repetition. Though it will come up again repeatedly, it is well here to dwell on the question who it is who speaks here. \randnotiz{Who is the Lord thy God?} Who is the Lord thy God? Dwell on His personality; on His unity in Trinity; on His holiness and hatred of sin;
\quellseite{217}
that He desires to be our Father, but can be a Father, in the full and comforting sense of the word, to those only who have the spirit of obedient children; that none can be such children unless they have been redeemed and renewed, are conscious of and grateful for this and desire to show this gratitude by loving obedience. Question on each point. Show that to those who willfully refuse to have God as a Father, He must be a Judge.

\randnotiz{Idolatry.} 2. Show what is meant by idolatry, in its gross and in its subtle sense. And how that anything or person that we fear, love or trust in more than in God, becomes our idol; show the danger and guilt of this sin and how it needs to be repented of, prayed over and guarded against. Illustrate and question.

\randnotiz{Fear.} 3. Show the meaning and duty of fear; that a God who has such wonderful attributes must be feared. That all mankind, even without a Bible, believe in some kind of a God or gods, and stand in awe of them; distinguish clearly between a slavish and a filial fear; show what kind of people have the one and what kind the other; show how each catechumen can know and ought to know whether he has a filial fear, e.\,g., whether he reminds himself frequently of God’s presence and knowledge of the heart and word and deed, whether he sincerely endeavors to guard against displeasing God, whether
\quellseite{218}
he grieves and repents when he has displeased Him, whether he would rather lose the favor of everyone else than of God. Here, already, you can catechise into the heart and experience. Be careful in your questions.

\randnotiz{Love.} 4. Dwell on the nature of love. You cannot define, but you can illustrate. Show the beauty and privilege of love; show why God is worthy of love; how love behaves itself towards God, delights to think of Him, to read His Word, to study it in the Catechism, to pray in private and to worship in Church, and how this love of God moves us to love our parents and our neighbor; warn against a false love, which loves the benefits and gifts of God rather than Himself, and speaks of love while the heart is cold. Here, again, you can speak to and appeal to the inner experience of the child. Question tenderly and cautiously.

\randnotiz{Trust.} 5. Show what trust means. Give examples; show that it is a thing of the heart, the feeling; show why God is worthy of unhesitating trust, warn against a wavering trust, which trusts only as far as it can see safety and benefit, but doubts in time of darkness and danger; show the blessedness of trusting always, in dark days as in bright days, in sickness and in health, in poverty and in plenty, in disappointment and in success. Impress and encourage here, for this trusting is one of the hard
\quellseite{219}
lessons for a Christian to learn. It is never learned perfectly here. Be careful not to promise too much with your illustrations. Show that in earthly things God often sees fit to disappoint and afflict for our good, but in spiritual things He always gives all we can or will receive. Question closely on each point.

\section{On the Second Commandment}

In every lesson the Catechist should ask himself seriously: What is the special sin, danger or error against which I should specially teach and guard my class to-day? This Commandment sets before us several specially frightful, yet common sins. One is the sin of cursing, the other of wrong or false swearing, and another is a superstitious resorting to would-be wonder workers, seers and revealers of what God has wisely hidden.

\randnotiz{Cursing.} 1. Here show (a) The uselessness of cursing. (b) Its coarseness and utter unworthiness of a lady or gentleman. Give examples. (c) Its awful sinfulness. A sin without even seeming advantage. An insult to God. Speak of an insult to mother. The only command with a threat.

\randnotiz{Swearing.} 2. Then speak of swearing; the trifling and unnecessary oath, the false oath, the awful sin of perjury; show how the perjurer calls down a curse upon himself; show when an oath is proper and how
\quellseite{220}
it should then be taken; show wrong of secret society oath.

(a) It is not demanded by the State, which alone has the right to demand. (b) The swearer promises under oath to keep secrets whose nature or content he does not yet know. This is unworthy of a rational and much more of a spiritual being. (c) The Masonic oath is a vile and hideous one. Show the beauty and heroism of unswerving truthfulness at all times and under all circumstances. Question always, even when not here indicated.

\randnotiz{Conjuring.} 3. As to resorting to the pretenders to the occult and hidden; show how it indicates a lack of trust in God and positive unbelief; how God has been good in hiding the future; how many of us could never have another happy hour if we knew the future; how irreverent and insulting to God it is to try to pry into what He has so wisely and so kindly kept concealed. Show the beauty of trusting and loving God, of doing our part and leaving the future with Him, of resting peacefully in His promised love and care.

\randnotiz{Duty or privilege.} 4. Then show the positive duty or privilege. How glad we ought to be to call upon God in every time of need, and how regularly we should do so; how we should worship, in the family, in the closet, in church, how in all we should pray, praise and give thanks; and that he who does this regularly and
\quellseite{221}
from the heart will always abhor cursing, wrong swearing, deceiving by His name, in resorting to divine healers, charmers, fortune tellers, spiritualists and other sacrilegious pretenders. Question closely and impress the need and duty of private prayer and church going and reverence in church. We should have said above that wanton irreverence in church, jesting or making light of holy things, are different ways of taking God’s name in vain. In bringing in illustrations and examples always draw on the Scriptures and especially on the life and example of Christ.

\section{On the Third Commandment}

\randnotiz{Holy day of rest—age.} 1. Show how the law of one holy day of rest goes back to creation, is hinted at in Noah’s dove and in the worship of the patriarchs, anticipates the Ten Commandments in the giving and directions on gathering the manna, and is taken for granted in the word, “remember.”

\randnotiz{For man’s good.} 2. That man needs such a regulation and that it is given for his good; that its purpose is rest for the body and refreshment for the soul; that whatever interferes with this purpose is wrong; that to go to church rests the body and refreshes the soul; that in proportion as we read, hear, teach and assimilate God’s Word, in that proportion do we hallow the day and get rich blessing out of it.

\quellseite{222}
\randnotiz{Failure.} 3. Show that there might be an outward abstinence from work or amusement, and yet a failure to keep the day, because the Word is neglected.

\randnotiz{Profanation.} 4. That, on the other hand, mere church-going and spending the rest of the day in sinful pleasure, in bad company, in unnecessary work, is also a profanation of the day.

\randnotiz{The Lord’s day.} 5. Show that to insist on the exact seventh day is ceremonial, legalistic, Jewish. But to rightly keep the Lord’s day, which is recognized by our risen Lord, by the Apostles, by the glorified Christ, and by the early Church, is evangelical loyalty to the Lord of the Sabbath; that the good Christian keeps this Commandment, as he keeps all the others as well as he can, not because he must, but because he wants to.

6. Guard your class against a puritanic rigorism on the one hand and against a worldly, pleasure-loving, continental holiday on the other. Question fully, and here as everywhere invite and encourage questions. All works of mercy and real necessity are allowed.

\section{The Fourth Commandment}

1. Show the high and honorable place that God has given to parents; that they are next to God; that children owe everything to them.\randnotiz{Why honor parents?} Remind them what parents did for them in infancy, in having them baptized, in childhood, in sickness, in sending
\quellseite{223}
them to school and providing for them, in giving them the privilege of catechisation—denied to so many.

2. Show how disrespect and disobedience are such fearful sins, that they are likened to apostasy and crime. Show what awful threats God’s Word utters against this sin.

\randnotiz{Superiors.} 3. Show that superiors in institutions, in schools, in church and in State, are in their sphere in the place of parents, and how our Catechism rightly demands for them what is due to parents.

4. Show what is meant by despising and by displeasing; how it is done; how shameful and sinful it is, and how it reflects on and avenges itself on the wicked doer. Give examples. Question.

\randnotiz{Duty to parents.} 5. Show how we should honor, how serve, how obey, how love, and how esteem them. Be very concrete and plain in your explanation and illustration of these terms. Show how it is to last for life, and what a disgrace and crime it is for children to leave their aged parents unprovided for, to let them suffer or go to the poorhouse.

Show that if parents should command what God has forbidden, then we must obey God rather than man.

As disrespect and disobedience are sins that are especially rife in our land and time the Catechist should take special pains to drive home and clinch the demands of this law.

\quellseite{224}
\randnotiz{Promise.} 6. Explain clearly the special promise attached to this Commandment. Make it clear that the promise is general; that as such it is true; that it means good parents; that children who follow the counsels of good parents do escape a thousand dangers and pitfalls into which the disobedient stumble and perish; that where good children of good parents are taken away by an early death it is for a special purpose and for a greater good than they could have by living. Never promise too much.

\section{The Fifth Commandment}

This one is easier to explain, but not easier to keep.

\randnotiz{Human life—sacred.} 1. Show sacredness of human life. God alone can give; He alone may take; He made man in His own image; who takes man’s life lays hold on and destroys the image of God. The awful sin of murder.

2. Show the even greater sin of self-murder. For this sin, committed by a rationally responsible person, there can be no repentance, hence no forgiveness. Warn most earnestly against this prevalent sin, you may save one.

\randnotiz{State and taking life.} 3. Show that God has delegated His right to take life to the properly constituted officers of the State; that, therefore, capital punishment is not murder, has God’s sanction, and ought not to be abolished.

\quellseite{225}
4. Show that God has sanctioned righteous war; that, therefore, to kill in such a war is not murder.

\randnotiz{Hate.} 5. Show in what sense hating one’s brother is murder; dwell on this; that God looks on the heart; that the wish is in God’s sight counted as the deed, that if not hindered through fear of consequences, it would often lead to the actual deed. Warn against spite, anger, envy and all bitterness of heart.

\randnotiz{Care of body.} 6. Show how, without committing actual suicide, one may so neglect, expose and abuse his body as to be guilty of bringing on suffering and decay; show that nature’s laws, health laws, are God’s laws; that nature never forgives a debt, but collects it to the last farthing. Take care of these precious bodies.

\randnotiz{Soul murder.} 7. In like manner show how one may endanger, poison and ruin one’s own soul and the soul of another, by neglecting and causing others to neglect what belongs to salvation, and by exposure to the blight and poison of sin. This is soul murder. Warn against it. Be very plain. Illustrate and question closely on every point.

8. Show how greedy employers and corporations are guilty when they provide no safeguards, overwork and underpay, and how by carelessness and thoughtlessness loss of life may be brought on.

9. Show how each one can and should assist and comfort in danger and want. Show beauty of an
\quellseite{226}
\randnotiz{Positive teachings.} unselfish life of benevolence. Illustrate from ministry of mercy and ministry of Word. Give concrete and encouraging examples, and show beauty and richness of such a life, in contrast with a selfish, passionate, peevish and bitter-spirited life.

\section{The Sixth Commandment}

We come now to the most delicate and difficult part of the whole Catechism for the conscientious Catechist with a mixed class. Of the advisability of dividing the sexes here we have spoken.

\randnotiz{No division of class.} 1. We do not favor division. We have shown what we hold to be a more excellent way. For warning against self-abuse, which is vitally necessary, he may invite the boys to spend an evening at his house. As to the girls, let him use his own judgment.

\randnotiz{So-called purity books.} 2. As to the recommending of books which minutely describe the sexual organs and their functions, together with warnings and hints against abuse and for healthy development, we do not, with our present convictions, favor them. For mature persons, with pure minds and established characters, they may serve a useful purpose. But for catechetical boys and girls, we fear that such books will suggest and teach more impurity and vice than safeguard against them. This at least is our present conviction.

3. Let the Catechist, however, speak with all possible
\quellseite{227}
gentleness and impressiveness. If he has won the love and esteem of his class, with this tone he can be very plain without provoking an impure thought or a smile.

\randnotiz{Warnings.} 4. Let him show the awful sin and the frightful consequences of impurity in thought, in word and in deed. Let him graphically paint the hideous picture of the defiled mind, the foul mouth, the filth of an impure life. Let him plainly tell of the blasted and blighted characters, the execrated society, the loathsome diseases, the dreadful early deaths and the despairing suicides that are the earthly wages of this sin. And then the eternal loss of the polluted soul. Not many questions here except such sermonic questions which are to be answered silently before God. For teaching this Commandment let there be special prayerful preparation.

5. Then let the Catechist reverse the picture. \randnotiz{Purity.} Portray the beauty of a pure heart, pure conversation and a pure life. Exemplify from Joseph, Job, and above all, from the stainless life of Jesus.

6. Let him point out and warn against the things that breed and lead to impurity, \emph{e.\,g.}, idleness, lying in bed after getting awake in the morning, intemperance in eating or drinking, a lazy, voluptuous life, impure books, shows and pictures, bad company, etc.

7. Let him point out and strongly admonish to the means that further purity, \emph{e.\,g.}, to be always
\quellseite{228}
occupied with something useful and good, to live abstemiously, to avoid stimulating foods and drinks, to avoid all the above dangers, to seek always the companionship of good books, especially the Bible, and of good people; to keep the example and suffering of Christ in mind, to never forget the all-seeing eye of God, to constantly watch and pray, realizing the danger of even the purest.

And then comes the great subject of marriage.

8. Here let him know that marriage is honorable, well pleasing to God, ordained by Him, and is the natural and Christian mode of life. \randnotiz{Marriage.} Let him show the sin, the danger and the sad individual consequences of the modern disparagement of marriage and motherhood; that home-life is the true life, and that it must come to ruin where marriage is despised. Speaking earnestly, he can speak plainly.

\randnotiz{Engagement.} 9. Let him speak most earnestly and urgently on the subject of engagement. Show that a proper engagement is the essence of marriage; that to choose a life companion is to choose one’s weal or woe; that the bonds of matrimony should draw to heaven, but may drag to hell; that to choose wrongly often means a life of martyrdom and an eternity of woe.

\emph{Ergo,} warn against hasty and ill-advised engagements; show what ought to be physical, mental, moral and spiritual bars to engagement. You cannot
\quellseite{229}
be too plain or urgent here, you may prevent some martyrdoms. Show how the counsel of parents should be sought, but that parents have no right to insist on marriage for social and pecuniary reasons where there is no mutual love, or where there are other objections that conflict with God’s idea. Earnestly admonish caution, delay, self-examination, acquaintance of the person and antecedents, earnest prayer. Question closely.

\randnotiz{Married life.} 10. Then as to married life. Hold up the ideal, love, bearing, forbearing, forgiving and constant devout reading of the Word, watchfulness and prayer.

Warn against spite-holding, jealousy, unbecoming and dangerous familiarity with the opposite sex, which often lead to impurity, separation and divorce.

\randnotiz{Divorce.} 11. Speak plainly of the calamity and criminality of the divorce evil; show what are not Scriptural grounds for divorce and what are; show the awful sin and baleful blight of divorce. Such catechising God will own and bless to the good of the class, the safeguarding of society, the stability of the State and the joy of the Catechist.

\section{The Seventh Commandment}

After the sacredness of the family and home we have that of property, without which the family could not exist.

\quellseite{230}
1. Over against the varied socialistic theories of the day let the Catechist here ground his pupils in the truth that God has from of old recognized and regulated the rights of private property, both real and personal.

\randnotiz{Property—acquirement.} 2. Show how property, both real and personal, can be acquired, viz., by gift, by inheritance and most generally by labor; that all other ways are wrong, \emph{e.\,g.}, cheating, taking advantage of another’s ignorance or necessity, such as by foreclosing an almost paid up mortgage, etc.; that all gambling, hazardous speculation, dealing in futures, cornering markets, and so making thousands of poor pay to make one or a few rich; the oppression of labor to enrich the capitalist; all combinations in restraint of lawful trade, are unrighteous ways of getting property. But let the Catechist be careful that he do not teach untenable, disloyal and unwarranted doctrines. He must be clear and sound in his own sociology and economics.

3. Show that much poverty comes from laziness, shiftlessness, intemperance and wastefulness, and in so far is a sin.

\randnotiz{Counsels.} 4. Show that to be neither too rich nor too poor, to have food and raiment and a good home, with a good education, a healthy body and a Christian character, is enough.

5. Show the right use of property, its stewardship
\quellseite{231}
and the responsibility for benevolence that it imposes. Here inculcate liberality, with what God has first given, toward the Church, her interests, institutions and operations. On all points question closely.

\randnotiz{Theft.} 6. Show how the Commandment is violated in robbery and theft. Warn children against the dangerous habit of taking or keeping little things that do not belong to them, \emph{e.\,g.}, spending the Sunday school change, etc.; by unfair dealing and fraudulent means, \emph{e.\,g.}, refusing to pay debts, taxes, stealing an employer’s time, car fares, idling away time in school and college.

\randnotiz{Duty to neighbor.} 7. Show how to assist our neighbor to improve and protect his own, by words of good advice, by helping where we can, by being kind and helpful always.

\randnotiz{To self.} 8. Show how one can wrong and rob himself, by not making the best of his opportunities, by not having a righteous ambition to make the most of himself, by ruining his health, by not guarding it, intemperance, or over-abstemiousness, by anything that injures body, mind or character, directly or indirectly. Illustrate and question.

\section{The Eighth Commandment}

As the former Commandment would guard the property, this one would guard the character and
\quellseite{232}
reputation. Luther’s Explanation is singularly full, clear, admirable and happy.

\randnotiz{Truthfulness.} 1. Show the importance of truthfulness; that without it confidence between men would be gone, that law and order would be at an end, that commerce could not continue, that society could not hold together, that the world would no longer be a fit place to live in. Exemplify and question.

2. Show the importance and preciousness of a good name, that without it life is a failure and happiness impossible. Illustrate and question.

3. Show that false witness is every use of the tongue, or of sign language, or of the pen that is against truthfulness and against true and unselfish love toward our neighbor. Illustrate and question.

\randnotiz{Sins against truth.} 4. Show how easy and how common and how ugly and how dangerous and sinful it is to fall into habits of deceitfulness, lying and spreading of evil reports; show the contemptibleness and the havoc of character, of peace and of all that sweetens life, of an evil tongue; show how lying and malicious gossip divide families, churches and neighborhoods, and leave only wreck and ruin in their train; urge the catechumens to examine themselves whether they are or would be guilty of such awful and inexcusable sins. Question closely, but so that they must confess and answer to God.

5. Show that fearless and absolute truth is specially
\quellseite{233}
important when bearing witness in court; the awful sin of perjury; show that when the court for the protection of society, or of an individual demands it, we not only may, but must tell what bad things we know against our neighbor. Make this very plain and show that the motive must be pure.

Question.

6. Make clear that a lie is an intention to deceive. Warn against actions intended to deceive, against being a hypocrite, against untrue exaggeration, false flattery, boasting, lies of politeness, etc. Illustrate and question.

\randnotiz{Positive duties.} 7. Show clearly what it means to apologize for our neighbor, speak well of him and put the most charitable construction on all his actions. Illustrate, exemplify. Show beauty and satisfaction of and commend these virtues most heartily. You can make the hearts, the lives, the homes, the congregations, the social circles and communities touched by your catechumens, brighter, sweeter, better by rightly catechising on the Eighth Commandment.

\section{The Ninth and Tenth Commandments}

We prefer to take these two together. They both treat of coveting. We believe many of the mechanical and artificial differentiations are distinctions without a difference.

1. Show clearly what coveting means; that it
\quellseite{234}
\randnotiz{Covetousness.} arises from the inborn inclination to evil; that it is original sin; that it rebels against God and the things and lot He has assigned us; that it is discontent, fault-finding, fretting, envying others, becoming bitter against them, and at last scheming and plotting to get what they have. Illustrate and question.

2. Show that for His own wise reasons and purposes God has neither made nor placed people alike. Some are stronger and brighter than others. It is better so. There is no dull monotony, yet all can be happy; some cups are smaller than others, yet each cup can be full; some are richer than others. What a dead and purposeless world if all were even! It is better so; it is not wrong, but laudable for each one to make the most of himself and to better his material condition. But to idly repine and rebel and become bitter, envious and covetous, not only does no good, darkens one's own heart and life, but is rebellion and sin against a wise and good God. It is a root sin. If unchecked it will lead to gross outward sin. Exemplify and question closely.

3. Show that in these last two Commandments God desires to open up to us our danger and warn us, by pointing us to the source of corruption within us; that He desires us to realize that His law is not satisfied with outward acts, but demands purity and integrity within. Question.

\quellseite{235}
\randnotiz{Covetous desires.} 4. Show that here God would have us guard against and repent of all evil and covetous desires; would have us be content with such things as we can rightfully have; would show us how we need to watch and pray against discontent, envy, bitterness and effort to obtain what belongs to others; would have us set our affections on the things that are above, seek first God’s kingdom and righteousness, be thankful that we can be sons and daughters and heirs of the heavenly King, rejoice in the prosperity and promotion of our neighbor and be ready always to further and help him, to assist and serve him in the preservation of his own. Question.

5. By these Commandments God would make us realize even more forcibly than in the others, that we are verily guilty of all, that with a heart so unsound by nature no one could or did keep the Commandments; that the law can never save us; that ourselves can never save us; that the law must drive us to Christ for salvation. Question and review.

\section[The Conclusion]{The Conclusion of the Ten Commandments}

\randnotiz{God and sin.} 1. Repeat and reimpress that God is holy; by nature abhors sin; must claim from us hearty abhorrence of and repentance for sin; wants our whole heart; our undivided and supremest love; is jealous and offended if we turn away and set the affections which belong to Him on other persons and
\quellseite{236}
things; He cannot be indifferent to sin, but must visit and punish sin, even to the third and fourth generation.

2. Here enlarge and enforce again how all just governments punish; how all good parents punish; how a just God must punish; show again how He punishes even in this life. Review on Commandments. Here is the place to counteract and fortify against that popular, present-day heresy that represents God as such an indulgent Father, that He shuts His eye to sin, takes even the willful and impenitent sinner to His bosom, and has no hell for the children of the devil.

\randnotiz{God and mercy.} 3. Show that while from His very nature He must punish—He could be not a holy God, a just God, if He did not punish—yet He delights in showing mercy. He shows it without stint or limit to thousands of them that love Him and keep His Commandments.

\randnotiz{Keeping the Commandments.} 4. Here show in what sense we can keep the Commandments, viz.: (a) Never perfectly. (b) Never in our own strength. (c) Never in such a way as to merit salvation by them. But: (a) Through the power of the Holy Spirit we can desire and strive to keep them and become more and more complete in our obedience. (b) This we do not to be saved by our obedience, but to show our gratitude for the redemption purchased by Christ
\quellseite{237}
and bestowed on us in our baptism and applied in His Word as we study it in these lessons. (c) And yet the more we strive to render a perfect obedience and the nearer we come, the more do we realize our remaining and besetting sin, and the more do we realize how we need to repent and fly to the dear Saviour for refuge and help.

Question and apply closely on every point. Rest not until you have made the most earnest effort to make each one see and feel his sin and repent.

To this end a verbal and Socratic review, with such searching questions and admonitions as this will give occasion for, will be helpful and blessed.

\quellseite{238}
\chapter{Chapter XXI. The Apostles’ Creed}

In the introductory lesson:

\randnotiz{Creed—meaning and history.} 1. Make clear the meaning of Creed. Show why the Bible will not do. Show need of creeds, use of, for Church’s comfort, for distinguishing truth from error. Show that everyone has a creed, if not written, then unwritten. It is important to impress and emphasize the need and use of creeds in our age which is hostile to them.

2. Give a brief history of this Creed. Its sacredness, because of its age and associations.

3. Distinguish between ecumenical or general and particular creeds. Give brief account of Lutheran Creeds. Have children memorize their names. Review and question on these three points.

\randnotiz{Creed’s relation to the Law.} 4. Make clear the relation of the Apostles’ Creed to the Law or Ten Commandments. Show the demands. It demands being right and doing right. Remind that through the fall man became so diseased and demoralized that he can never fulfill the demands of the law. He must be saved in some other way. The Creed, especially in its second and third articles, shows the way. Question.

5. Insist and impress that neither mere memorizing,
\quellseite{239}
nor mere mental acceptance as truth is enough. Dwell on the three elements of faith, viz., \randnotiz{Elements of faith.} knowledge, assent and trust. Distinguish carefully. Show that trust is a thing of the heart, an experience. Question closely on this whole point. It is one of the points that ought to grip the heart and conscience. Show that only a heart-acceptance fits for Confirmation.

\randnotiz{Three articles.} 6. Show relation of the three articles to each other: The first treats of Creation. Because man, created in God’s image, became sinful and corrupt, Redemption was necessary, and this purchased Redemption needs to be applied in renewal or Sanctification.

\randnotiz{The Trinity.} 7. Make clear the doctrine of Trinity. God reveals Himself as One, also as Three, Three present at Christ’s baptism, Three in baptismal commission, yet one name. Show that being unable to understand and grasp, is no reason for doubting or disbelieving. Illustrate from other mysteries in Scripture, in nature, in daily life. Question.

\randnotiz{The Trinity in salvation.} 8. Show how our creed sets forth that each Person has a share in the work of our salvation; that our hearts’ experiences require and demand the three; sometimes more a Father, as when in danger and fear, again a Redeemer, when our sins trouble us and seem too great to be forgiven, and still again, when we realize and mourn a spiritual dearth and
\quellseite{240}
coldness, a quickening, reviving, refreshing and strengthening Spirit. Question closely and fortify pupils on the doctrine and the use of the doctrine. It is unwelcome and unpopular in our day, and our catechumens will be sorely pressed for believing it.

\randnotiz{Attributes.} 9. Here it might be well to bring in the principal attributes of God. Be careful to define in simple Saxon words. Steer clear of the terminology of dogmatics.

Explain each attribute until you can question the idea out of the child. Show the use and application and comfort of each attribute. Clearness here is worth more for the average child than abstract and scientific speculation on the being and nature of God.

\section{The First Article}

\randnotiz{God—nature and existence.} 1. In this age of naturalism and Darwinian evolution it is well to dwell on the words, “I believe in God.” The Catechist must fortify his catechumens on every point where they are likely to be attacked and shaken in their faith. Our age does not believe in our personal God. If Catechist can set forth and lodge the simple sense, in simple words, of the common proofs for the existence of God, here is the place to do it. But if he cannot do it without using such terms as “cosmological,” “ontological,” etc., it is better that he say nothing, for such learned terminology will only confuse the child. He is not a
\quellseite{241}
good Catechist if he cannot so give the substance of each proof that he can draw intelligent answers from the children.

Then show again that a mere mental assent is not enough, but that this tremendous truth of the existence and presence of that God who has the attributes set forth above, must move the heart and stir the conscience. This cannot be insisted on too often.

\randnotiz{Creation.} 2. Next, speak of God’s creation. Drill them on the creation story until they know and can repeat the successive order. Do not insist on the twenty-four hour day. Show the word day is used in various senses in Scripture. But be careful not to go too far and do not raise questions which you cannot answer.

3. Here tell of the creation of angels, and of the fall of those who became bad angels or devils; of the existence of a real, organized kingdom of darkness, with Satan as its head, who does not have the divine attributes of omnipresence, or omniscience, but is informed and works through the legions of subordinate spirits. Fortify on this point; show clearly the Scripture teaching as to a personal devil and his angels; show how he tempts, endangers and degrades his victims till they become children of the devil; how this doctrine explains the frightfully inhuman deeds of bad people, and so helps to relieve the mind and explains what cannot be explained if
\quellseite{242}
there be no devil and his angels. Question. Then tell of the good angels, that they are our ministering spirits, serving and protecting God’s children. Exemplify from Scripture and question.

\randnotiz{Our Father.} 4. Dwell on the fact that God is our Father, the Father of all by creation and preservation, but the Father of those who have been redeemed by Christ and renewed by the Holy Spirit in a special and comforting sense. Fortify and caution against the false rationalistic and popular idea of the Fatherhood of God, without any heart acceptance of Christ as a personal Redeemer. Question.

\randnotiz{Preservation.} 5. Show what God does for each one as a wise and good Father: Gives and preserves body and soul. Show that body was made out of earth, can be satisfied with what is in the bowels of, \emph{e.\,g.}, metals, money, stones, fuel, etc., and what grows on surface, or from earth-nourished animals, or other earthy human bodies.

That soul was not made of earth, but was breathed in by God, is related to God and can be satisfied only in God and things of God. Make plain and telling application of this wonderful dichotomy, which is taught in this confessional book, and show the worth of and the sin of neglecting or abusing the God-given and preserved body, soul, reason, faculties, raiment, food, home, family, property. Dwell on, apply and question on each
\quellseite{243}
word. Do the same with God’s daily providing and protecting. All this shows a rich and kind Father.

\randnotiz{Why.} 6. Show why He does it: Without any merit; show that we can merit nothing but punishment, because we are such sinful children; show that it is all goodness and mercy. Dwell on, apply and draw out these points.

\randnotiz{Obligations.} 7. Show what each one owes for all this. Dwell on each word. Show that outward thanks and praise are nothing without willing and obedient service; that in this also God wants the heart; that such serving and obeying is not a burden, but a joy.

Question.

8. Show that each one can have and should have the unshaken conviction that all this is most certainly true, that in addition to the proofs already set forth each one’s own heart and conscience testify that this is all true. What manner of persons, then, should we be, if all this is true?

\section{The Second Article}

Here is the most beautiful, the most precious part of the Catechism. It would be hard to find any uninspired writing equal to Luther’s explanation of the second article. Here is the heart of all evangelical theology, the foundation of all true Christian ethics. The article is too rich and full for us here to exhaust it, even proximately. Happy is the Catechist
\quellseite{244}
who daily experiences its preciousness and knows how to explain, apply and make precious its reach treasures of truth.

We can give only hints.

1. Dwell on and explain and question on the names of our Lord given here and in other parts of the Bible; show how each name has instruction and comfort for God’s children. Question.

\randnotiz{Names and offices.} 2. Do the same with His offices of Prophet, Priest and King. Dwell on the Prophet. Explain and apply. Show diligently and earnestly what we should learn from this divine Prophet. Question. Explain only the names and functions of Priest and King, as the priestly and kingly work are brought out more fully in the further explanation. Be sure that the children are clear on each step.

3. Show again why it was necessary to have a Redeemer. Why a mere man could not redeem us. Why God could not, without becoming man. Make clear and question.

\randnotiz{True God.} 4. Show why we believe that Christ is true God. Here the Catechist must lay himself out and be strong. No doctrine of the Word and Church has been so violently, so skillfully and so insidiously attacked as the divinity of Christ. Every Catechist needs to strengthen himself by a repeated study of Dr. Whitteker’s “Separated Life.” Drill the catechumens till they can answer every objection and
\quellseite{245}
defend every statement. Make it specially clear that Jesus claimed to be God; that if He made a false claim, He is a deceiver; that He must be God, or He is an impostor. Thus show that there is no standing-ground for the sweet-mouthed rationalist and Unitarian. Rub this in and draw it out. Marshal and drill in the divine names, attributes, works and worship. Question on each point. Show the satisfying comfort of the true doctrine, and vice versa.

5. Show that He is also true man. \randnotiz{True man} Human birth, human body, human soul, grew up like another child in body, had human wants and emotions, wept, suffered, died. Therefore I have a Saviour who can understand me, feel for and with me, is the very kind whom I can approach without fear and is able to redeem and help because He is God and man in one person. Question closely.

6. This Christ redeemed me. Dwell on the beautiful and significant meaning of the word redeem, show comfort in it.

7. Me—what kind of a me?—a lost and condemned creature. \randnotiz{Lost and condemned} Remind how I became lost, and why I was condemned. Dwell again and impress again the unwelcome doctrine of sin, sinful nature, guilt and condemnation. Call Scripture, history and observation to witness. Dwell on infant sinfulness, on infant suffering and death, which is sin’s wages.
\quellseite{246}
Call their own conscience to witness. Question closely. Surely I needed to be redeemed. Emphasize the “I” and “me” and show the significance; show that Luther’s Catechism is a handbook of experimental religion, as well as of doctrine.

8. \randnotiz{From what redeemed} From what has He redeemed me? Sin, death, devil. Show in what sense I am redeemed from sin, viz., from its guilt and from its power. In what sense from death? That though I must still die, the sting of sin is gone, all fear of death is gone, death is an entrance to the real, the eternal life. In what sense from the devil? Where sin is forgiven the devil has neither claim nor power. Because I am redeemed I can resist him and make him flee. Question and apply.

9. \randnotiz{With what redeemed} With what has He redeemed me? Not, as slaves were redeemed, with silver and gold. His whole earthly life of humiliation, birth, suffering, death, burial, all this was vicarious. In each step He was my substitute. Dwell on vicarious atonement; show that Gethsemane’s agony and the despairing cross cry, “My God, my God, why hast Thou forsaken me?” are meaningless and inexplicable if the atonement was not vicarious. This is another disputed, disowned and hated doctrine, Ergo, take special pains here. Make the true doctrine so clear, and fortify it so forcibly and so fully, that all can see; state and uphold it. It really has
\quellseite{247}
fewer difficulties than any of the other theories of the atonement that have been brought forward to supplant it. Question and cross-question and repeat.

10. \randnotiz{His exaltation} Show how after the atonement was finished by His humiliation, the whole work of redemption was divinely sealed, proclaimed, acknowledged and made sure from heaven by His triumphant descent into hell, His glorious resurrection, ascension and session at God’s right hand—which is everywhere. And that all will be finally crowned when He comes again to judge the living and the dead. Dwell on and explain each point. Show what each point means to each one personally and individually, and how each one should daily remind himself of these wonderful truths and their relation to himself. Be sure to bring in and apply each Scripture, which we do not indicate in these hints, but which the Catechist is to scrupulously select, explain and apply at every point. Question as on last point.

11. For what has Christ thus wonderfully redeemed me?

paragraph

\randnotiz{Why redeemed} That I might be His. Emphasize the “I” and “His.” I belong not to myself nor to the world. He has paid for me the incomprehensible price of that vicarious atonement. I cost Him much. More than any tongue can tell. Impress this. Tell of Zinzendorf’s view of the crucifix with the inscription: “This have I done for thee, what hast thou done for me?” Question.
\quellseite{248}
He wants me to belong to Him, to yield Him my heart, my love, my life. He has a kingdom for me. Grace here. Glory there. I am the son, the daughter of a King. What a privilege! What endowments!

And serve. Serve Him. I am redeemed not merely that I might keep out of hell, not that I might selfishly enjoy the fruits of my redemption, not that I should be satisfied that I get to heaven. No, no! Such a selfish religion is unworthy of one redeemed at such a cost. Such selfishness is a shame and a disgrace. Question.

I am redeemed to serve; to serve Him; to serve in righteousness—a holy service. In innocence—a pure and upright service. In blessedness—a glad, a happy service. Everlastingly—an endless service. Dwell on this. Be concrete. Show opportunities for service. Endeavor to win candidates for the blessed service of the ministry, the ministry of mercy, or the ministry of the Word. Review and reimpress and question on the whole article and show again the blessedness of one’s saying it from a warm and loving heart.

\section{The Third Article}

\randnotiz{Who the Holy Ghost is.} 1. Make clear who the Holy Ghost is. Not an influence, not an impersonal power, but a person in the same real sense in which the Father and the Son are persons.
\quellseite{249}
The third person of the Trinity. True God, names, attributes, works, worship. Drill on this. Much uncertainty and doubt are abroad. Question closely.

\randnotiz{His work.} 2. Explain the work of the Holy Spirit. That He takes the things of Christ and brings them to us, i. e., He applies what Christ purchased. His work, therefore, is to renew and sanctify. Question.

\randnotiz{Order of operations.} 3. Show the order of these operations. Luther does not mention baptism here because it comes up fully in the fifth part. He speaks here of the Spirit’s operations in those who are old enough to read, hear and understand. These He calls, gathers, enlightens, forgives or justifies, sanctifies and preserves, and finally raises from the dead and glorifies. Explain each word clearly, until the children can give their own meaning. Distinguish each operation from the others and show their relation. Question.

\randnotiz{His means.} 4. Show by what means He does this, viz., the gospel and His gifts. The gospel here means the whole Word, including the law. The gifts signify any other means and instrumentalities, e. g., my instruction in these things, God’s providential dealings with me. Other gifts are the Sacraments, which will be treated separately.

Here dwell specially, pointedly and plainly on the Word as the Spirit’s vehicle or means through which
\quellseite{250}
\randnotiz{The Word} He brings His renewing and sanctifying grace into my soul. Show clearly how the Word, including law and gospel, the written Word, when read, when preached, when pondered, when connected with the Sacraments and so becoming the Sacramental Word, how this Word carries the Spirit, is His instrument, and as such does His blessed work. Enlarge on and emphasize our Church’s doctrine of grace through the means of grace, in which doctrine she stands alone, which is one of her chief treasures and glories, which bridges the chasm between God’s sovereignty and man’s responsibility, which is so full of light and comfort; ergo, not my own reason or strength, but the Spirit through the means. Rest not till your pupils understand that the Spirit operates through the Word, and that we have no right to expect His holy influence and power apart from the Word. Question closely on this important and helpful point.

\randnotiz{The Church} 5. Show where this Spirit operates, viz., in the Christian Church. Here show what the Church is, visible as the congregation of believers among whom the gospel is purely preached and the Sacraments rightly administered; invisible wherever there are true children of God, known only to God; show that Christ builded the Church, loved her and purchased her with His blood and committed to her the means of grace; show that the Church is called
\quellseite{251}
His bride, the Lamb’s wife, and that in the end He will take her home and celebrate the marriage supper of the Lamb.

It is important to dwell on and emphasize this point, because in our day ministers, and even bishops—not of our Church—preach that one can be as good and safe outside as in the Church. The Church is held in low esteem, and multitudes prefer human organizations and lodges to the Church of Christ. Question closely.

\randnotiz{The purest Church.} 6. Here it will also be in place to remind your class that the Church is unhappily divided, that there are many so-called churches. Which is the right one? Without going into subtle refinements which your catechumens cannot grasp or appreciate, it is well to impress this one self-evident truth, that that is the best and the purest church which sticks most closely, believes most heartily, confesses most fully and fearlessly and practices most consistently all that God has clearly revealed in His Word. By this test our own Church is willing to be tried and judged. Other churches may be good on some points of Scripture doctrine, ours comes the nearest to being good and loyal on all Scripture doctrine. Make this very plain to your class and be not satisfied until all can give satisfactory answers on this point. Show them here and frequently what they would lose if they should desert their own Church for another.

\quellseite{252}
\randnotiz{The Last Things} 7. Take up the doctrine of the Last Things, viz., the resurrection of the body and the life everlasting. Here also our age is skeptical. Fortify your catechumens that they make not shipwreck of their faith and lose this blessed hope and comfort.

\randnotiz{Richness of third article} 8. On this whole third article an old and pious Lutheran saint has said:

“This whole article is so rich in the most important matters that can engage the heart of man. These things should ever be so explained and impressed that the hearts of the catechumens will be moved to allow the Holy Spirit to carry on His gracious work in each of them. Should there, however, be a Catechist who does not himself stand in the order of grace, and who has not experienced these gracious operations of the Holy Spirit in his own heart, such an one will be sadly cold and careless, even while he hurries over these living and life-giving lessons.”

May God grant that this may never fit a Catechist who reads these pages.

\quellseite{253}
\chapter{Chapter XXII. The Lord’s Prayer}

1. The catechumen now understands how and through what means the Holy Spirit applies Christ’s purchased redemption, renews and sanctifies. In so far as he has taken this to heart, is conscious of having been personally called, gathered, enlightened, forgiven, or justified, and is in the process of being sanctified; in so far as he experiences these gracious operations of the Spirit, in so far is he a consciously trustful child, yearning for closer communion with God, and able to cry, Abba, Father! \randnotiz{When prayer is real} In so far he can now really pray. Question closely as to who can really and truly pray.

2. Make clear the nature and conditions of acceptable prayer.

\randnotiz{True prayer} True prayer is the trustful conversation of a true child with his heavenly Father. God must first speak to me, make Himself known to me, and be understood by me. I must realize that He speaks to, assures, comforts and promises me all that is good for me. When I have thus heard, heeded, taken to heart His assurance, comfort and promise, then, resting on that Word, holding up to my loving Father His own Word and promise, basing my confession,
\quellseite{254}
petition and thanksgiving on that Word, I truly pray.

\randnotiz{The Word and prayer.} Question closely on the relation between the Bible, as God’s Word to me, and prayer, as my word to God. Show clearly and draw out that the Word must always be studied prayerfully, and that prayer must always be combined with devout study and contemplation of the Word.

Show that in the Word God comes to, offers to and bestows on me, while in prayer I bring to and ask from God. Ergo, the Word is sacramental, while prayer is sacrificial. Ergo, prayer is not a means of grace. Illustrate. Asking for bread does not appease hunger, but bread does. Question.

\randnotiz{Answer to prayer.} 3. On conditions of acceptable prayer make clear the distinction between spiritual or eternal good on the one hand and bodily or temporal good, including the health and life of the body, on the other; that when we really long and pray for the former it is always God’s will to hear and answer us, and that the answer is in proportion to the heart’s desire. But when we pray for the latter, i. e., for earthly good, we remember that God is the wise Father, we the ignorant children. We do not always know what is good for us. He does. Illustrate from relation between earthly father and child. Paul’s thorn in the flesh. Question closely on distinction between heavenly and earthly good, on relation between father and child.
\quellseite{255}
Show that therefore God does not promise to hear every prayer, and therefore when we pray for earthly good our prayers must always be conditioned and say, “Thy will be done.” Show that when God withholds one desired good He gives another and better, e. g., Paul’s thorn. Question closely and make clear these distinctions. Show wrong and right in faith-healing.

\randnotiz{Prayer and God’s order} 4. Show that we have no right to expect God to set aside His own means or methods of grace. If the Spirit works through the Word, I have no right to pray God to convert a person without the Word. I must see that the means are brought to that person and then pray God to bless those means. Illustrate and question.

Even in physical disease I must use means or remedies. God’s Word sanctions physicians and physical remedies, ergo, again, Christian science is wrong, and fanatical faith-healing is wrong. Question closely. In all this introduction bear in mind that many misleading and hurtful ideas on prayer are abroad, that the Reformed idea is different from the Lutheran; that it is the Catechist’s duty to get his catechumens clear on all these points.

\randnotiz{Model prayer} 5. After making clear the nature and conditions of true prayer, show that this, the Lord’s, is the best of all prayers, and a model for all true prayer. It is the best, because of Him whose it is and who
\quellseite{256}
commends it, because it embraces all our wants in a few words, because of the beautiful order of the petitions.

\randnotiz{Our Father} The word Father reminds me of my relation, out of this child and Father relation all true prayer must flow. The last word, Amen, means faith, without which there can be no true prayer.

Again, “Father” is a word of trust; “our” is a word of love; “in heaven” is a word of hope. Faith, love and hope are here. Question on all this.

\randnotiz{Order of petitions} As to the order of the petitions: The first three deal with our spiritual needs, the middle one with our temporal needs, and the last three with spiritual deliverances. To hallow His name should be the ground and aim of all our spiritual strivings, this is done in His kingdom, where His will is supreme, and the doing of it should be our soul food and delight. Then we pray for bread for the body. In the getting and use of it we sin much against God and against our neighbor, are sinned against by him, ergo, we need forgiveness and need to forgive; and we are in constant danger of falling again. Ergo, we pray against temptation, and against the roots of sin or evil. Question.

6. In all this, show how this prayer is a model for our prayers. Impress that the catechumens are now old enough and ought to make their private prayers in their own words, telling God their special
\quellseite{257}
\randnotiz{Prayer—one’s own.} and particular needs and dangers and praying for the special help they need. And impress deeply the need of a prayerful spirit by all who really desire to be and remain true children of God. Impress the need of regular prayer, e. g., morning and evening, especially before retiring. Be not too legalistic on posture. Show that the proper private posture is kneeling, but that on a really cold night in a cold room, if they can pray better in bed, God will accept the prayer, if only it is real and sincere.

Show also that in addition to their regular prayers, all should form the habit of sending up silent petitions in every time of special need and danger. Show that reading or repeating from memory does not detract from the true worth of prayer, if only the heart is in every petition, and that making one’s own prayer may be pure formalism if the heart is not in it. On all this question generally, but be careful not to tempt to a hypocritical or deceitful answer. Your great aim here is first to give clear ideas on prayer and then to induce to the habit and life of prayer.

7. On the introductory words, follow Luther’s explanation. Show again what is implied in the endearing name Father,\randnotiz{The introductory words.} and what it implies on the part of each child. Dwell again on the child relation, its privilege and blessing. Show what it means to call on Him with all cheerfulness. Who can do
\quellseite{258}
this? With all confidence, who can do this? Show the altruism in Our. In heaven, reminds of heavenly attributes. Dwell on and review and apply again each attribute. He is on earth, also, omnipresent. If my Father is such, what should I, His child, be? Question closely and practically.

Show again the falsity of the popular conception of God’s universal Fatherhood, and the conception of an indulgent father, who makes light of and overlooks sin. Warn against this false idea.

Impress again that only he who has Christ as a personal Redeemer, and lives in personal, trustful and loving relation with Him, can really love God as his Father. This is important. Question closely. Remind of fear, love and trust.

\section{The First Petition}

1. Review Second Commandment and its explanation. Show its relation to this petition, showing holiness of God’s name; show that hallow means to recognize as holy and to treat as holy, especially in private and public worship. Question fully.

Show when this is accomplished, show meaning of teaching, expounding, preaching the Word in its truth and purity, i. e., to be careful and prayerful always to understand correctly and to accept completely. \randnotiz{Importance of pure doctrine} Here dwell on and impress the importance of pure doctrine, the irreverence, danger and sin of
\quellseite{259}
caring not whether we understand God aright or not, the danger of present-day laxity, looseness and liberalism; show how it casts a slight on God and charges Him with insincerity in His revelation, that a liberal church or preacher is irreverent, unsafe and unworthy of our esteem. This is important. Be clear and impressive, but not polemical. Be sad rather than angry at the irreverent indifference to truth that is abroad. After you have made the danger of carelessness clear and shown the scripturalness of loving sound doctrine, question closely.

2. Show that by having and following a few important principles of hermeneutics or explanation we can know what the teaching of the Word in truth and purity is.\footnote{See “Way of Salvation in the Lutheran Church,” pp. 33, 34.} Drill and draw out these hermeneutical principles.

\randnotiz{Sound doctrine of Lutheran Church} 3. Here is a place to emphasize (a) that our Lutheran Church loves sound teaching, has made the most thorough research and investigation of what the Word teaches on every point of doctrine, and (b) that it is a privilege and comfort to belong to such a sound Church, and (c) that it would be a most serious loss to leave such a careful and safe Church for one that is indifferent and unsound on some of the precious teachings of the divine Word. Impress these points and draw them out.

\quellseite{260}
4. After having emphasized the sound doctrine it is of vital importance to emphasize even more earnestly the demand of a holy life.

\randnotiz{Holy life demanded.} Impress that sound doctrine alone is not enough, that one might have this and be and remain a stranger to Christ and His grace, and be eventually lost. Impress and draw out this momentous truth. Then show what a holy life is, how it must harmonize with the Word and will of God. Review the third use of the law. Show that it demands holiness in thought, in feeling, in purpose, in word, in act. Emphasize that only he who earnestly aims and strives after such holiness is a good Lutheran Christian. For grace to grow in such holiness of heart and life we must daily pray that our heavenly Father may help us. Question and apply.

5. This petition also prays for help to be kept from profaning the name of God. Who does profane that name? He who teaches otherwise? Again the awfulness of false teaching. And lives otherwise. Again the awful danger of unholy living. And even more so if there is a pretense to holiness, while the heart is cold and the life unholy. Warn and question. Impress the need of watchfulness and prayer.

\section{Thy Kingdom Come}

1. Distinguish and define the three phases of
\quellseite{261}
\randnotiz{God’s kingdom—three phases} God’s kingdom, viz., the kingdom of power, of grace, of glory; show that we must be in the second before we can be in the third; show that God has always had His kingdom of grace and of glory; that the kingdom comes indeed without our prayer. But we pray that it may come to us. We should never come into the kingdom of glory if the kingdom of grace did not come to us. Grace comes to us through the means of grace. We pray that the Holy Spirit may enable us so to use the means that they may bring the kingdom into us.

But we should pray that the kingdom of grace may come to all men, to the heathen who are without the means of grace; we pray here for missions. Luther came short in not including the missionary idea in his explanation. If we claim to be Christians and are unwilling to help the heathen, we are untrue to our name. True children of God pray and give and labor that the kingdom may come to the heathen also. Question, illustrate, apply.

\randnotiz{Kingdom through the Spirit} 2. Show that we cannot have the kingdom within us unless our heavenly Father gives us His Holy Spirit. Review explanation of third article of Creed. Therefore the need to pray that God would give us His Holy Spirit. Question.

3. If we have that Spirit the kingdom has come to us. Then we believe His Word. We believe it all. We believe even where we cannot understand.
\quellseite{262}
But here also our faith is not a dead faith, not a mere intellectual assent, not mere historic belief, but a faith that warms the heart, that works, that enables and impels us to lead a godly life on earth. This again describes a true Lutheran Christian. Then we shall finally pass from the kingdom of grace into the kingdom of glory and keep on living this godly life in heaven forever. Question.

\randnotiz{Kingdom of darkness} 4. It might be well to refer here to the kingdom of darkness with the devil at its head. Show that it is an organized kingdom, pitted against the kingdom of God, and that all who are unwilling to be in the kingdom of God must be in this kingdom of darkness in this world and in the next. Question. Review again on devil and bad angels under Creation in first article of Creed.

\section{Thy Will be Done}

\randnotiz{God’s will} 1. Show that God’s will is a good and a gracious will; that it is done without our prayer, but that we pray here that it may be done by us also. Enlarge, illustrate and question.

2. Show how the will of the devil, the will of the world, i. e., of the unbelieving world, and the will of our own flesh, i. e., of unregenerate human nature and its remnants in the best of us hinder us and prevent His kingdom from coming to us. Enlarge, illustrate
\quellseite{263}
and question. Review on devil and his angels.

3. Show how God frustrates and brings to naught the evil counsels and purposes of these enemies, \emph{e.\,g.}, \randnotiz{God overruling.} by His restraining, limiting and punishing providence, by bringing good out of intended evil, by disciplining and strengthening His children through these testings. Exemplify and question.

4. Show how we need to pray that God may keep us steadfast in His word, and in our faith, that we are weak, cannot keep ourselves, but He can and will.

5. Show how perfectly and joyously God’s will is done in heaven, that we should strive to approach that glad and trustful obedience more and more; \randnotiz{Doing God’s will as in heaven.} that in order to do this we need to discipline ourselves (a) to acknowledge that God’s will is good and gracious even when we cannot understand His permissions and doings, and even when they seem to be sorely against us; (b) that we bow under the chastening rod, knowing that what God inflicts is good, even when it hurts; (c) that we do not blame God for the inflictions and hurts we bring upon ourselves by our own disobedience and sin. Distinguish clearly here; show how we make ourselves suffer when we disobey parents and the laws of health, and how wrong it is to call such sufferings mysterious providences. Be clear yourself on providence
\quellseite{264}
and take pains to make your catechumens clear. Unclearness here is responsible for much skepticism. Question closely.

\section{Our Daily Bread}

1. Show what is embraced under bread, \emph{viz.}, every earthly good, as Luther exemplifies. Show that these are all good in themselves; how that indirectly they are all God’s good gifts. \randnotiz{Meaning of bread.} Review explanation of first article. Distinguish between necessities, comforts and luxuries, between enough and superfluity; show that we must often be content with necessaries; have a right to pray and strive for comforts, but have no right to covet luxuries and envy those who have them. Review Ninth and Tenth Commandments. Dwell on these points. Here lie the roots of an unchristian socialism. Question closely.

2. Show how God gives, \emph{viz.}, by giving strength and opportunity for labor, by gifts, by inheritance. Review Seventh Commandment. Show that God promises nothing to the lazy. Bread is promised for labor, thrift. \randnotiz{How God gives.} Show again sin of wastefulness of God’s gifts. Impress and question.

3. Show how God gives even to the wicked; that His goodness should lead them to repentance; that there is still a great difference between their having and that of God’s children; that the former are
\quellseite{265}
selfish, thankless, without inner appreciation and joy; that the better we are as God’s children, the more content and joy we have in God’s gifts. Remind of former petition as to God’s will if we are providentially poor. Be very plain and practical on this petition. Question closely. Invite questions always and everywhere, but especially here. Endeavor earnestly to allay doubt and discontent and to satisfy sincere misgivings as to God’s dealings.

4. Show in what sense daily bread or earthly good is ours; not bread obtained by theft or fraud, but bread and all earthly good honestly obtained and in that sense ours. \randnotiz{When bread is ours.}

5. Show that we need it daily, should pray for it daily, and, doing our part, should trust God for the next day. Remind of explanation of first article. Question clearly.

6. Show that our bread, or food, is God’s gift; it should always be received with thanksgiving. Table prayer, always, either aloud or silent, \emph{e.\,g.}, in a public place, but always heartfelt. Warn against formalism, saying without thinking or meaning, and so taking God’s name in vain. \randnotiz{Thanksgiving.} Question generally. But do not invite deceit.

7. Do not forget to call attention to the “us.” It is not “my.” It is not selfish. It means concern and interest in the needs of others.

\quellseite{266}
\section{Forgive us, as we Forgive}

1. Our debts, our trespasses, our sins, sins against God, sins against our neighbor; show again the sinfulness and damnableness of sin. Review this point as it came up in conclusion of Commandments and in each article of Creed. Be not afraid of repetition. The doctrine of sin is a root doctrine, \randnotiz{Sinfulness of sin.} an unwelcome doctrine, a disputed doctrine, a hated doctrine. Fortify your class. Question closely.

2. Show again that all have sin, that there can be no sinless perfection, that all need to pray for forgiveness. \randnotiz{All need forgiveness.} Question.

3. Show the wretchedness of an impenitent and unforgiven heart, the blessedness of forgiveness, the conditions, true penitence and faith. Explain again. Show that forgiveness is justification; show again that sin had to be atoned for before it could be forgiven, and how Christ did atone; that His atonement was vicarious; that, therefore, God can forgive and delights to forgive, where the heart is in condition to receive it. \randnotiz{Explanation of forgiveness.} Question closely.

4. Show how we sin against our neighbor by neglecting to assist and comfort and protect. Remind of explanation of Fourth, Fifth and Seventh Commandments. Show how we sin against him in our thoughts by envying and hating; in our words by false and evil speaking, in our deeds by injuring his person or his property. Remind of Command- \randnotiz{Sin against neighbor.}
\quellseite{267}
ments that forbid all this, especially of the Eighth and its beautiful explanation. Question.

5. Dwell again on the sin of bitterness, of spite-holding, of its disastrous consequences to self, to family, to church, to community; show that it is unworthy of a Christian. Make this clear and drive it home and question closely.

6. Then make clear that a bitter and unforgiving heart cannot be a broken and a contrite heart, and that God cannot forgive that heart; that that heart prays, “Forgive me, as I forgive, but I forgive not at all, therefore forgive me not at all.” Show the frightfulness of this petition uttered by one whose heart is bitter and spiteful against another. Question closely. \randnotiz{Unforgiving not forgiven.}

7. Show that we must forgive as God forgives; that while He demands penitence and confession, He wearies not in coming to the offender, reasoning with him, pleading with him, beseeching him to be reconciled. And so must we do with those who offend against us. Make this also clear and emphatic and question closely.

8. God not only forgives, but He pours out His richest blessings on those whom He forgives. We therefore should be specially kind to those whom we forgive and always readily do good to them all. \randnotiz{As God forgives.} Exemplify and question. Show the beauty, the blessedness and the happiness of thus forgiving and
\quellseite{268}
returning good for evil. Give examples and question.

9. Show again the constant need of earnest prayer for our own forgiveness and for a forgiving and loving spirit.

\section{Temptation}

1. Define temptation clearly. Make catechumens memorize definition. Show what temptation to evil is, that such temptations come from the devil, the world and our own flesh. Give examples. Show awful sin of belonging to the world and so tempting others to sin. Give examples and question. \randnotiz{Temptation defined.}

2. Show that all are in danger. Devil is at work through under-devils and devilish people who are in his service. Examples, question. Show again reality of personal devil. Review what was said on bad angels under first article. Tell again his story from Bible, beginning with serpent in Eden; show how a wicked world helps him along, and how our own flesh is inclined to yield. \emph{Ergo}, all are in great danger and need to pray earnestly, lest they fall. \randnotiz{Danger.}

3. Show that God cannot tempt to evil. Distinguish clearly such temptation to evil from testing, trying, disciplining anyone for his own good; show that in this sense God tempts, tests, tries and disciplines for the strengthening of faith and character. Exemplify and question. \randnotiz{God tempts not.}

\quellseite{269}
4. Show that we here pray for strength to overcome the temptations to evil, for God’s help that our enemies may not succeed in deceiving us, leading us into error, \emph{i.\,e.}, into false belief, \emph{e.\,g.}, that we are strong and safe and there is no danger; that there is no devil and no danger from one; that God is indulgent and overlooks sins, or that God is not good, else we would not suffer and be tempted; show danger of such errors or false beliefs, and safety in true doctrine. Show danger of being led in unbelief as to God’s mercy and Christ’s atonement, and the Holy Spirit’s help, and how such unbelief may end in despair and spiritual death. \randnotiz{Why we pray thus.}

5. Show that on account of the weakness of our flesh, and the strength of the devil and the world, we need to watch and pray, lest we fall into other great and shameful sins. Examples and close questions. \randnotiz{Why watch and pray.}

6. Show how that if we are in earnest in this petition we will never expose ourselves to the danger of bad company, bad places, where the influence is dangerous to spiritual life. Review Sixth, Ninth and Tenth Commandments and Providence under first article and above petitions. Show that no one who needlessly exposes himself to danger has a right to expect God’s protection. Question.

7. Show that each one should be earnest and diligent in the use of the means of grace, that so he
\quellseite{270}
may grow stronger and ever stronger against temptation and sin.

\section{But Deliver us from Evil}

1. Show that this petition is a summary. Dwell on every word of Luther’s explanation. Question on every point.

2. Show again how all evil comes from sin. Where there is no sin, there is no evil, no hardship, no suffering of any kind. The root evil is sin. For deliverance from it we here pray. If we have deliverance from sin there is no danger for the soul; bodily evils can be cheerfully borne, while poverty, losses and disappointments can also be met with Christian fortitude and courage. If sin is forgiven and its power broken, there is comfort in the midst of all the ills that may afflict our earthly life. Impress this, your catechumens will need it. Question closely. \randnotiz{Evil from sin.}

3. Show again how the earthly ills and crosses which we do not bring upon ourselves, but which God lays on us, are all meant for our discipline, for our training, our growth in grace, our weaning from the love of earth, the setting of our affections on heavenly things, the transforming of us more and more into the image of Christ. Make this clear. \randnotiz{God’s crosses.}

4. Show that when we pray this petition we also
\quellseite{271}
pray that we may be made more and more willing to bear whatever crosses God may impose and as long as He may impose them, that we may not become rebellious or wish to end our life, before God says it is enough. Question.

\section{The Conclusion}

1. Show that in this conclusion the child of God shows his reason for believing in and coming to God in prayer. Because the kingdom and the power and the glory are His, therefore I know that He can do for me whatever He deems good for me. To such a God I can pray. What a privilege to know Him, to be permitted to come to Him with all confidence and cheerfulness, even as a beloved child to an earthly parent. Question. \randnotiz{Reason for prayer.}

2. Impress every word of Luther’s beautiful explanation and application of the word Amen. \randnotiz{Amen.} I should be assured, made certain, beyond all doubt, that such petitions are acceptable, are heard. Blessed assurance! What an encouragement to pray. Impress it. Show the emptiness, the comfortlessness, the fearfulness, the hopelessness of a life without prayer! Drive this thought home, so that they may never forget it.

He has commanded prayer, in this manner, after this model. May I ever study it, pray it, and conform all my own prayers to it. He has promised
\quellseite{272}
to hear, i. e., when I pray after this manner, according to the foregoing explanations, in the spirit inculcated above. Then He hears. Amen, Amen! Yea, yea, it shall be so!

3. Show how easy and how frequently this precious prayer is abused. It is so well known, so frequently used, that there is great danger of saying it without praying it, of rattling it off without either thought or heart in it. This is a sin. It is taking God’s name in vain. Remind of Second Commandment. Warn earnestly against this sin. Luther once called this prayer the greatest martyr on earth, because it is so much abused. Question and caution. \randnotiz{Abuse of prayer.}

4. Remind again that to pray aright, and especially to pray this prayer aright requires (a) That I see and grasp the mind, the meaning, the will of Christ in it. (b) That I have an earnest, heartfelt longing to have and use every good I pray for in each separate petition. (c) That I may have the spiritual wisdom to apply each separate petition to my present need. Impress and draw out these three points. \randnotiz{Praying aright.}

\quellseite{273}
\chapter{Chapter XXIII. The Sacrament of Baptism}

We have repeatedly admonished the Catechist that he must know the danger points in the catechumen, the points on which he is in danger of faltering and falling. Our last three chapters deal with such points. In her teaching on Baptism, on Confession and on the Lord’s Supper the Lutheran Church stands alone. All the Reformed churches differ from the Lutheran on these points. The Lutheran who imagines and tries to show that there is little or no difference here is either strangely ignorant or wilfully indifferent and disloyal to what his Church has always taught, confessed and defended against those who hold and teach a different doctrine. Let the Catechist be clear here. It is a serious matter to betray the teaching of the Church whose name he bears. \randnotiz{Lutheran Church position.} It is a grave responsibility to teach error to the unsuspecting catechumens. It is a woeful wrong to them to leave them unclear, unfortified, an easy prey to plausible errorists, who will find, attack and trouble them and possibly rob them of true comfort for the rest of their life. “Take heed unto thyself and unto the doctrine.” “Be not carried about with divers and strange doctrines.”
\quellseite{274}
“Hold fast the form of sound words.” “Earnestly contend for the faith.” “Be ready always to give an answer.” “The time will come when they will not endure sound doctrine.” Is not that time here? Are not our catechumens in danger?

1. Remind and impress again that our Church is the only Church that has a clear and consistent teaching on the fact that grace comes through means. That where the means are, there grace is. That we know of no grace and have no right to expect grace where there are no means. \randnotiz{Grace through means.}

2. Make clear again that the Word and the two Sacraments are the means. Show again that the Word is the principal means, that it makes the Sacraments means. Show again that in other churches neither Word nor Sacrament are believed to be means that carry, convey, bring into the heart and life the grace that bringeth salvation. Read carefully chapter seventeen of “Way of Salvation.” Then chapters three and four, then thirteen, fourteen and fifteen. Question closely on what constitutes a means of grace. Show again the difference between the Lutheran view and the Reformed view of the Word of God as the bearer of the Spirit and grace of God. \randnotiz{The means of grace.}

3. Give and impress a clear definition of a Sacrament. Show what three marks or characteristics it
\quellseite{275}
must have. Show that only Baptism and the Lord’s Supper have these three marks; that Rome is wrong in claiming seven Sacraments; show types of the two in Old Testament; show that without a believing and loving heart they are of no avail. That depending on the mere outward observance of our Sacraments, while the heart is cold, faithless and loveless, is not merely useless, but seriously sinful. Question closely. \randnotiz{Sacrament defined.}

4. Before catechising on Baptism read carefully “The Lutheran Pastor,” pages 296-299. Examine yourself as to whether you are prepared to safeguard and fortify your catechumens against the specious attacks that they will certainly have to meet from the baptistic and revivalistic sects. Many Catechists fail here and are largely responsible for the future apostasy of their catechumens.

5. Make clear the nature of baptism. Make clear the meaning and bearing of Luther’s first two questions. Emphasize that it is not a ceremony gotten up by man, or by the Church, but that it comes from the hands of the God-man, Christ Jesus; that He would not institute and make binding an empty or useless ceremony; that if it were man’s we might ask, what is the use? or what good can a little water do? But it is dishonoring to Christ to ask such questions about an institution of His. He can do great things through little means. Illustrate and impress these fundamental points. \randnotiz{Nature of baptism.}

\quellseite{276}
6. The Catechist must be clear on the question of the baptism of infants. He must be clear on the doctrine of original sin. He must realize that the Lutheran alone has a sure ground for infant baptism because of his Scriptural views of sin, of grace and of grace through the means of grace. He must know what the arguments against these teachings are and have an answer, a satisfying answer, for every one. \randnotiz{Infant baptism.}

7. Being clear on the ground and need of infant baptism, he must know all the historical and Scriptural arguments that prove the Lutheran position correct. Dr. Seiss’ “Baptist System Examined” is still one of the most satisfactory books on the subject. Whitaker’s little book on Baptism is popular and helpful.

8. Now, let him review the children again on sin, on the need of renewing grace, on the means, on the nature of baptism, the Scripture proof, infant baptism and invalidity of arguments against. Drill catechumens until they can answer every argument. On Baptist challenge to show an example of infant baptism in Bible, show, first, that God established infant membership in what Stephen called the church in the wilderness, that God never recalled it, that no one else dare; second, that there were infants, hosts of them, among those who crossed the Red Sea, that Paul says they were baptized, that \randnotiz{Infant baptism—Scriptural.}
\quellseite{277}
when Paul calls it a baptism we had better not contradict. Third, show that the commission to disciple the nations, by baptizing and teaching, includes infants, as there never was a nation without infants. Illustrate from census taker. Next show that there is not a single example in New Testament of the Baptist custom, \emph{i.\,e.}, of a child born of believing parents and baptized as an adult; that all the adult-baptisms recorded in New Testament are those of converts from Judaism or from heathenism. Then show that the history of the early Church shows that infant baptism was the custom by all except by those who were accounted heretics. Drill patiently and perseveringly on every point. Time and patience used here may save some from going astray.

9. The less important matter of the mode of baptism may be brought in next. This is more important in some communities than in others. It is not likely, however, that all the catechumens will escape future onslaughts from fanatics who have more zeal than knowledge. The books referred to above will show the Catechist what the opponents’ arguments are and how he can meet them. Let him treat this subject in the manner recommended for baptism of infants above. Drill in the arguments against immersion and for sprinkling and in the answers to the objections of opponents. Impress upon them that they were proper subjects when they were \randnotiz{Mode of baptism.}
\quellseite{278}
baptized and that the Sacrament was properly administered to them. And impress further that since this is most certainly true, it will be disloyalty to Christ, to His Church and to His truth, if they ever repudiate their baptism and receive one of another kind. These things cannot be made too clear or too strong. Right teaching here may save shipwreck of the faith hereafter.

10. Then come Luther’s two questions on the benefits of baptism and the promises of God.

Here make clear again the condition in which the child is without baptism. Then again show how God the Father loved those children and gave His Son to redeem them, to purchase for them the forgiveness of all their sin; how the Son effected this purchase of pardon when He became their substitute, took their sin upon Himself, endured the guilt and suffered the punishment which the sin deserved, and so made forgiveness possible. But this purchased pardon must be personally and individually applied. This God the Holy Ghost effects through the water and word of baptism. The Triune God has thus entered into covenant relationship with His baptized children. Their sins are forgiven, death and the devil can harm them no more—everlasting salvation is theirs. They are the sons and daughters of God. Thus should the Catechist impress upon the children the priceless treasure they have in their \randnotiz{Benefits of baptism.}
\quellseite{279}
baptism. It has brought and sealed to each one individually all the blessings of salvation. Show how the faith that was implanted potentially or germinally through baptism must be nurtured and maintained, that only they who keep on believing what the Word and promise declare shall be saved.

And here show the children again how they can tell whether they still have the life and faith given in baptism. Make clear to them that if they can make light of sin, laugh at it, and say they don’t care, they have reason to fear they do not really and heartily believe. Show that with the heart man believeth, and that he whose heart has faith must hate sin, mourn over it and earnestly desire and seek forgiveness when overtaken by it. Impress the need of being in earnest about this deliverance from sin, from the fear of death, from the power of the devil. Impress the fact that children who are utterly thoughtless, careless and prayerless, and willfully continue so, have no right to comfort themselves with their baptism; that these comforts belong to him who really believes and continues to believe. \randnotiz{Baptism not appropriated.} Warn against a false comfort, as if the mere fact and form of baptism assure salvation. Emphasize the fact that in baptism God has made a covenant with each one, that He will be faithful; but so must the recipient be. If the baptized child wilfully breaks this covenant and becomes a prodigal,
\quellseite{280}
the covenant is broken. But it is not God’s fault. On His part the covenant stands, and He will receive the prodigal child back when he repents. All this must be made clear to the children.

11. Coming to Luther’s question: “How can water produce such great effects?” Remind again of the power of the Word as a means of grace, show that the Word is essential; that water is necessary because it is a part of Bible baptism. \randnotiz{How these great effects.} Remind of the place of faith, as above; show in what sense it is a washing of regeneration; what regeneration means; that it is elsewhere called being born again; that being born, or a birth, is a very feeble beginning of an independent life, and yet a real beginning of a real life; that baptism is the means, the gracious water of life, the start of the child of God. Show again the preciousness of this gift of God, and the base ingratitude of despising it and throwing it away.

12. In Luther’s last two questions explain clearly what is meant by the old Adam. What by daily sorrow and repentance. Emphasize the “daily”; show that the child of God never gets done repenting, because he never gets done sinning; \randnotiz{The old Adam.} that this daily repenting is a true sign, a certain manifestation of the new life. The old Adam, the old nature, has been deadened by the implanting of a new life. But that old nature still has some life. It shows itself in
\quellseite{281}
sins and evil lusts. These need the daily drowning by daily sorrow and repentance. No child of God gets beyond this. There is no sinless perfection. “Forgive us our trespasses” needs to be a daily prayer. For this daily conflict, this daily overcoming of the old Adam, the sins and evil lusts, we need the law, we need the gospel, we need to watch, to pray, to wrestle. Impress the earnestness of the Christian life, the constant danger and the constant striving against evil. Question closely and searchingly. Let no one pass this point without feeling its reality.

Show how the new man is to come forth and rise. That as through baptism the old man is crucified with Christ, so the new man is risen with Christ. As the crucifying is to go on daily, so is the rising again to be a daily effort and exercise; that there is to be a renewal day by day. Our growing Christian life is thus to be a constant fellowship with His suffering as well as a living into the power of His resurrection. \randnotiz{The new man.}

Show that so living our baptism will be a daily and never-failing spring of comfort and of strength; that the grace from this means of grace will abide as a precious uplift through life. Question and impress and apply all this personally, until each one realizes: “This means me and is for me as if I were the only one.”

\quellseite{282}
13. Remind them again that our Church is the only one that has this rich, precious and Scriptural doctrine on baptism; that if any of them should leave this Church for another they would throw all this away and rob themselves of all this help. Stir up here also an appreciative love for the Church that teaches them so clearly, comforts them so effectively and enables them to hope so steadfastly. \randnotiz{Preciousness of Lutheran doctrine.}

14. Here remind them of the relation of their baptism to their coming Confirmation. Study over again the chapter on Confirmation and see that you make the meaning, the responsibility and the blessedness of this act clearly understood. Rest not until you have done all in your power to make clear, drive home and clinch the full significance of a voluntary Confirmation. And warn again against that superficial, state church and rationalistic view and habit of Confirmation which has wrought such havoc in our dear Church. Resolve, by the help of God, that your catechumens shall understand the real spiritual self-surrender of this beautiful rite, and that to each one Confirmation shall mean an intelligent and a voluntary consecration for life. \randnotiz{Baptism and Confirmation.}

\quellseite{283}
\chapter{Chapter XXIV. Of Confession}

Between parts four and five of the Catechism, as a bridge from one to the other, looking back to the grace vouchsafed in baptism and looking forward to a worthy reception of the Lord’s Supper, we have Luther’s three questions on confession.

This intermediate part of the Catechism brings before us a doctrine peculiar to our Church. She stands alone in her belief and teaching on this subject. Other Protestant churches have no confessional teaching on this subject, unless it be to deny and repudiate what we teach. In the Reformed churches our teaching is abhorred and denounced as Romish and dangerous to the soul. Some Lutherans also are afraid of it, avoid mention of it as much as possible, handle it very gingerly when called on to speak of it, and often pervert it. To the Reformed ear and taste it is one of the most unpopular of all church doctrines. \randnotiz{Peculiar to Lutheran Church.}

This aversion arises mainly from ignorance of the meaning and intent of the doctrine. It is prejudice. People who harbor the prejudice have never sincerely studied the subject. The very mention of confession and absolution seems to repel them.

\quellseite{284}
Unfortunately, many Lutheran ministers seem to be unclear on the subject. They seem to be unable to make it clear to others. Their teaching and preaching leave the impression that the minister can and does in his own name and person forgive sins. Their careless, impenitent and formalistic catechumens and hearers flatter themselves that because they have heard the pastor read the absolution, therefore their sins are ipso facto forgiven. Their misunderstanding of the doctrine makes them feel secure in their impenitence and sin and endangers their soul’s salvation. Here, then, is a blessed opportunity for the good Catechist. How important that he rightly explain and apply these catechetical questions on confession and absolution. We offer a few hints to help him. \randnotiz{Doctrine misunderstood.}

1. Read again carefully chapter sixteen of the “Way of Salvation.” Also chapter fourteen of “The Lutheran Pastor.”

2. Make clear the Romish teaching on the subject. Show what is false and unscriptural in it. Show falsity of the hierarchical conception of the priesthood as an order. Show impossibility of enumerating all sins, even if the priest could forgive them, the uneasiness and fear of the soul lest some sins might have been unconfessed, and, therefore, unforgiven. \randnotiz{Rome’s false doctrine.}

3. Show that there are three kinds of confession,
\quellseite{285}
viz., one to God, one to a fellow-man, when we have wronged him, and one before the pastor. Rehearse and emphasize again the daily, constant need of penitence and confession before God,\randnotiz{Kinds of confession.} not only on account of the known sins, but on account of the secret sins, the sins of the heart which are known to God alone. Question closely.

Then rehearse and reimpress the duty of confessing our faults one towards another. This passage is often misapplied. It has nothing to do with confessing to the pastor, unless one has wronged him personally. Show again the sin of holding spite, the special importance of personal reconciliation before partaking of the Holy Communion. Question personally, drive it home.

4. This brings us to Luther’s first question. Here show that it is not absolutely necessary to forgiveness that we confess before the pastor.\randnotiz{Confession before the pastor.} But that for the sake of the absolution and the comfort that it brings to weak and troubled hearts, such confession has been retained as a useful church custom.

Show that such private confession means a confidential, frank and free interview with the pastor, in which the heart and life are laid open, besetting and dangerous sins are fully confessed and comfort and counsel sought from one’s spiritual guide.

Show that such confession should be made with heartfelt sorrow, with a humble spirit, with earnest
\quellseite{286}
thoughtfulness, with child-like trust in our heavenly Father, and with earnest purpose to watch, pray and strive against such sin more earnestly in the future than in the past. Question on each point. Make each point clear.

5. Show that all sin is against God. That God alone can forgive sin. That no mere man, be he layman or minister, can forgive sins by his own power or in his own name. The pastor is God’s messenger, the mouthpiece and ambassador of God, when he brings God’s message, speaks God’s Word. That Word, as we have seen, carries God’s grace, carries forgiveness, life, salvation. In that sense the forgiveness announced and declared from the pastor is as from God Himself. It is not the pastor, but the Word, that bears the forgiveness. Recall illustration of government’s deputies, after Civil War, in chapter sixteen, “Way of Salvation.” \randnotiz{How sin is forgiven.}

If the Catechist can make this clear then all the Reformed objections fall to the ground, while all the comfort of the absolution remains. Thus rightly presented, this is another of the precious doctrines peculiar to our Church. Happy is the Catechist who can make it so clear and can so lodge it in the mind and heart of his catechumens that they can gladly explain and defend it against all gainsayers.

6. This explains Luther’s second question sufficiently. The third is so easy and practical that no
\quellseite{287}
more is necessary than that the Catechist impress each clause. Let him remind and reimpress the catechumen how to use each Commandment as it is needed; show them that while they are sons and daughters now, they will be servants, or masters, or fathers or mothers by and by; and how they will need this admonition of the Catechism as long as they live. Question on each point. \randnotiz{Luther’s third pedagogic use of the Ten Commandments.}

7. Show them how the lessons of this brief part of the Catechism are intended to help them to appreciate and get rich blessing out of the fifth and last part.

\quellseite{288}
\chapter{Chapter XXV. The Sacrament of the Altar}

In this, the last part of our Catechism, we again have to do with a subject on which our Church differs from all others. Here, again, it is of the utmost importance that our catechumens be made clear on our doctrine and on its grounds and proofs. As in the two preceding chapters, so in this. All the companions, friends and neighbors of our catechumens who belong to other churches will be against them. They will endeavor to persuade them that the Lutheran teaching is false, Romish, superstitious. All the non-Lutheran commentaries, religious books, Sunday school helps and other literature are against us. Our doctrine is unpopular. Lutherans themselves are often unclear, hesitant and skeptical here. The responsibility of the Catechist is great. He must first of all be clear himself. He needs to know the teaching of others. He needs to know and be ready to answer the objections to our doctrine. He dare not spare himself any pains to make this all so clear that it becomes a conviction, and to so equip his catechumens that they can answer the arguments of the opposers. \randnotiz{Clear statement of doctrine.}

1. Read carefully chapters thirteen, fourteen and
\quellseite{289}
fifteen of the “Way of Salvation.” Get its view-point, its argument, its illustrations. If a more comprehensive study is desired there is nothing in English more full or forcible than Krauth’s “Conservative Reformation.”

2. Explain clearly the different names of the Sacrament, e. g., the Lord’s Supper, in what sense a Testament, Sacrament of the Altar, Table of the Lord, Sacrament of the Body and Blood of Christ, Holy Communion. Show that each name indicates that we are here dealing with a sacred mystery. \randnotiz{Different names.}

3. Recall the time and circumstances of the institution, its testamentary character, sacredness of testamentary words, Paul’s inspired explanation of the words. Question. \randnotiz{Time and circumstances.}

4. On Luther’s first question show in what sense we have the true body and the true blood really present in the Sacrament; show that we have to do with the spiritual, glorified body and blood of the glorified Christ, who is everywhere present, has all power and can be specially present and give Himself to us in, with and under the bread and wine. Show that to doubt this possibility is to doubt the power and Godhead of Christ. Question.

5. Make very plain what the Romish doctrine of transubstantiation is, that according to it there is no more bread and no more wine, as the bread has been changed into flesh and the wine into blood, so that
\quellseite{290}
there is no more earthly element; show that this is against Paul’s explanation, who calls it bread after the consecration and calls it the communion of Christ’s body. Make clear that it takes two things to make a communion, and that both must be present. The Romish doctrine thus contradicts Scripture. Question fully. \randnotiz{Romish teaching.}

6. Show what the teaching of our Church is, viz., that the bread is not changed, neither is the wine. But these earthly elements, without change or material mixture, convey the body and blood, and that thus the bread is a real communion of Christ’s body and the wine of His blood. Many Lutheran Catechists are not clear here. They leave the impression that we believe in transubstantiation. They leave the same impression by their preaching. Question and drill. \randnotiz{Lutheran teaching.}

7. Show clearly what the Reformed doctrine is; that it gets rid of the heavenly element and retains only the earthly. It contradicts Paul’s teaching of a communion on the opposite side from Rome. Show how empty this teaching is and how it robs the Sacrament of all mystery and of all objective character, and really robs it of all sacramental character. Show that this loose doctrine leads to loose practice; that the Reformed people come without a preparatory service, that they invite everyone to partake who deems himself fit or moved to come; \randnotiz{Reformed teaching.}
\quellseite{291}
show that our doctrine cannot allow us to be so liberal and so loose. Question until you make this clear.

8. Now go over all the Scripture proof again, beginning with the words of institution and explaining them by the words of Paul, and show how our doctrine is loyal to the Word, and that it is dangerous to change the Word or to refuse to believe, because we cannot understand; show that it is better to trust too far, to believe too much, if this is possible, than to trust and believe too little. Question. \randnotiz{Scripture proof.}

9. Show the preciousness of our belief. How it makes this Sacrament blessed. How much the Sacrament means to us. How much it offers to us. How it becomes the Holy of Holies, the place nearest heaven this side of heaven. And what a comfort, what a joy, what a fuller view and value of our Saviour we have when we believe that He who gave Himself for us here gives Himself to us, as meat and drink for our souls. Show again that those who believe not our doctrine cannot have this comfort and joy in this Holy Sacrament. Question closely. \randnotiz{Our belief precious.}

10. Coming now to the benefits from the sacramental eating and drinking, show in what sense and manner forgiveness, life and salvation are granted; show that while the true child of God always has forgiveness, yet he needs the constant assurance
\quellseite{292}
from the Bible read and recalled, that the assurance comes with a deeper impress from the preaching of the gospel of forgiveness, that it comes home with still greater comfort in the words of absolution, but that it comes with the greatest of all comforts when the penitent believer takes the sacramental bread and wine from the pastor’s hand; he receives it personally, individually, as if he were the only one. With it he hears the words, “given and shed for you for the remission of sins.” He can no longer doubt that it means himself. As sure as he hears, takes, eats and drinks, so sure is he that his sins are all forgiven. With forgiveness he receives an increase and strength in his spiritual life. This increase of life is intended to increase his faith, love, hope, patience, watchfulness, prayerfulness and all the Christian virtues. Show that each one of these virtues ought to increase and grow with each communion. Show that thus salvation, from the guilt and from the power of sin, for the life that now is and for the life to come, is made more real and more precious by this holy Sacrament. Impress again that all this is only for him who has true, heartfelt penitence and a trusting, living, loving faith.

11. Show that the unworthy, i. e., the secretly impenitent and unbelieving, do indeed receive the same body and blood of Christ, but, instead of partaking of the above blessings they eat and drink damnation to themselves. \randnotiz{The unworthy.}
\quellseite{293}
Warn most earnestly and solicitously against such unworthy participation.

Review briefly the lesson on confession. Question to find whether all are clear on worthy and unworthy communing. This is a vital point. Serious mistakes are made on both sides. While some commune who are not fit, others fail to commune because they are unclear as to what real fitness is. They do not know that this is a means of grace not for those who are already good enough, but for those who realize their unworthiness and hunger and thirst after righteousness. Right catechising here will make fuller and more blessed communions hereafter.

12. Make clear your position on the frequency of Communion. Here also there are faults on both sides. But among us the most common fault is too much neglect of this blessed privilege. Read pages 345–347, “The Lutheran Pastor.” Warn especially against the serious sin, so frequent among our confirmed members, of communing after Confirmation and thereafter rarely, or never. Show that the inner life must be at a low ebb and in danger of expiring where there is no felt need and no desire for this heavenly food. \randnotiz{Frequency of Communion.}
